\documentclass[11pt, a4paper]{article}
\usepackage{amsmath}
\usepackage{amssymb}
\usepackage{graphicx}
\usepackage{hyperref}
\usepackage{listings}
\usepackage[margin=1in]{geometry}
\usepackage{fancyhdr}
\usepackage{booktabs}
\usepackage{tabularx}

\lstset{basicstyle=\ttfamily\small, breaklines=true, language=Python}

\title{ECG Simulator Documentation\\Part 8: Advanced Reference \& Future Roadmap}
\author{Extended Features, Phase 4+ Planning, Troubleshooting}
\date{June 2026}

\pagestyle{fancy}
\fancyhf{}
\rhead{Part 8: Advanced Reference}
\lhead{ECG Simulator}
\cfoot{\thepage}

\begin{document}

\maketitle
\tableofcontents
\newpage

\section{Troubleshooting Guide}

\subsection{Flat QRS (No Visible QRS Complex)}

\subsubsection{Symptoms}
\begin{itemize}
    \item ECG shows flat or minimal voltage during QRS interval
    \item P waves visible (normal)
    \item T waves visible (normal)
    \item Only QRS appears depressed or absent
\end{itemize}

\subsubsection{Root Causes}
\begin{enumerate}
    \item \textbf{Depol Gaussian decay}: QRS Gaussian has small $\sigma_d$ (default 12.7 ms)
        \begin{itemize}
            \item Solution: Increase depol\_duration to 40 ms for Phase 3 (LV overlap)
        \end{itemize}
    
    \item \textbf{Amplitude scaling too low}: amplitude\_scale < 1.0
        \begin{itemize}
            \item Solution: Set amplitude\_scale = 1.0 or higher
        \end{itemize}
    
    \item \textbf{Engine selection wrong}: Using Phase 1/2 with Phase 3 parameters
        \begin{itemize}
            \item Solution: Ensure engine selector matches code logic
        \end{itemize}
\end{enumerate}

\subsubsection{Resolution Steps}
\begin{enumerate}
    \item Check GUI amplitude\_scale slider (should be ~1.0)
    \item Verify Engine = "graph" (Phase 1/2) or "cable" (Phase 3)
    \item For Phase 3: Ensure depol\_duration parameter = 40 ms
    \item Click Reset button to clear old traces
    \item Click Start to generate new beat
\end{enumerate}

\subsection{No ECG Signal Generated}

\subsubsection{Symptoms}
\begin{itemize}
    \item Click Start button; no ECG trace appears
    \item Console shows no errors
    \item Display remains blank
\end{itemize}

\subsubsection{Root Causes}
\begin{enumerate}
    \item \textbf{Timer not running}: GUI timer not started
    \item \textbf{Parameters invalid}: Infinite or negative values
    \item \textbf{Engine disconnected}: Engine not properly initialized
\end{enumerate}

\subsubsection{Resolution}
\begin{enumerate}
    \item Check console output (enable debug\_mode = true)
    \item Verify HR is between 30--220 bpm
    \item Verify all conductances are 0.0--1.0
    \item Reset GUI: Click Reset button
    \item Try different engine (graph vs cable)
\end{enumerate}

\subsection{Irregular Rhythm When Expecting Regular}

\subsubsection{Causes}
\begin{enumerate}
    \item \textbf{AF enabled}: af\_enabled = true
        \begin{itemize}
            \item Solution: Set af\_enabled = false
        \end{itemize}
    
    \item \textbf{Ectopic focus active}: ectopic\_enabled = true
        \begin{itemize}
            \item Solution: Set ectopic\_enabled = false
        \end{itemize}
    
    \item \textbf{Escape pacemaker interfering}: escape\_enabled = true
        \begin{itemize}
            \item Solution: Ensure AV conductance = 1.0 so escapes don't fire
        \end{itemize}
\end{enumerate}

\section{Advanced Parameter Tuning}

\subsection{Matching Patient-Specific ECG Patterns}

To simulate a patient's ECG with specific morphology:

\begin{enumerate}
    \item \textbf{Measure patient ECG}:
    \begin{itemize}
        \item HR, PR, QRS, QT intervals
        \item QRS axis, QRS amplitude
        \item Presence of pathology (bundle block, ischemia, etc.)
    \end{itemize}
    
    \item \textbf{Adjust simulator parameters}:
    \begin{itemize}
        \item Set HR to patient HR
        \item Prolonged PR → increase av\_delay
        \item Broad QRS → reduce lbb\_conductance or rbb\_conductance
        \item Left axis deviation → enable LAHB (set lahb\_conductance → 0)
        \item Right axis deviation → enable LPHB (set lphb\_conductance → 0)
    \end{itemize}
    
    \item \textbf{Adjust timing constants} (Phase 3):
    \begin{itemize}
        \item Modify $\tau_s$ in cable\_1d.py to match patient's conduction times
        \item Re-run numerical search for phase calibration
    \end{itemize}
    
    \item \textbf{Validate against patient ECG}
    \begin{itemize}
        \item Compare intervals and amplitudes
        \item Iterate parameter adjustments
    \end{itemize}
\end{enumerate}

\section{Phase 4 Development Plan}

\subsection{Phase 4a: 3-D Ventricular Tissue (Q2 2027)}

\subsubsection{Objectives}
\begin{itemize}
    \item Extend Phase 3 cable model to 2-D ventricular sheet
    \item Implement orthotropic diffusion (longitudinal > transverse CV)
    \item Model 3-D dipole orientation changes during repolarization
\end{itemize}

\subsubsection{Technical Changes}
\begin{itemize}
    \item Replace 1-D Laplacian with 2-D: $\nabla^2 v = \frac{\partial^2 v}{\partial x^2} + \frac{\partial^2 v}{\partial y^2}$
    \item Introduce diffusion tensors: $D_L > D_T$ (anisotropy)
    \item Update CFL condition for 2-D: $\text{CFL} = D_{\max} \Delta t / \min((\Delta x)^2, (\Delta y)^2)$
\end{itemize}

\subsubsection{Expected Outcomes}
\begin{itemize}
    \item Improved QRS morphology (especially V1--V3 precordial leads)
    \item Realistic T-wave vector changes
    \item Ability to simulate infarct patterns (regional conduction delay)
\end{itemize}

\subsection{Phase 4b: Full Conduction System (Q3 2027)}

\subsubsection{Objectives}
\begin{itemize}
    \item Upgrade Phase 3 cable to full 3-D geometry
    \item Include:
    \begin{itemize}
        \item SA node (pacemaker dynamics)
        \item Complete atrial geometry
        \item AV node (two-dimensional with slow-fast regions)
        \item His bundle and bundle branches
        \item Complete LV and RV
    \end{itemize}
\end{itemize}

\subsubsection{New Parameters}
\begin{itemize}
    \item Fast-pathway conductance (AV node)
    \item Slow-pathway conductance (dual AV node pathways)
    \item Accessory pathway (Wolff--Parkinson--White syndrome)
\end{itemize}

\subsection{Phase 5: Advanced Cell Models (Q4 2027)}

\subsubsection{Courtemanche Model (Atrial)}
\begin{itemize}
    \item 21-variable state vector (vs 2 for FHN)
    \item Detailed ion channels: $I_{Na}$, $I_K$, $I_{Ca}$, $I_{K1}$, ...
    \item APD typically 100--150 ms (clinically accurate)
    \item Will enable: AF dynamics, APD alternans
\end{itemize}

\subsubsection{ten Tusscher Model (Ventricular)}
\begin{itemize}
    \item 19-variable state vector
    \item Epicardial, midmyocardial, endocardial variants
    \item APD typically 200--300 ms
    \item Will enable: Transmural dispersion, long QT syndrome
\end{itemize}

\subsection{Phase 6: Boundary-Element Forward Model (2028)}

\subsubsection{Replace Single-Dipole with BE Torso Model}
\begin{itemize}
    \item Construct 3-D torso geometry (from CT/MRI templates)
    \item Compute BEM gain matrix (heart → surface electrodes)
    \item Project 3-D ventricular dipoles to 12 standard + additional leads
\end{itemize}

\subsubsection{Expected Improvements}
\begin{itemize}
    \item Realistic lead vectors (St.--axis angles match clinical)
    \item Ability to model unusual electrode placements
    \item ST-segment morphology improvements
\end{itemize}

\subsection{Phase 7: Autonomic Nervous System (2028)}

\subsubsection{Sympathetic/Parasympathetic Modulation}
\begin{itemize}
    \item Model beta-adrenergic effects: increased HR, faster AV conduction
    \item Model parasympathetic effects: decreased HR, AV block risk
    \item Parameters: \textbf{sympathetic\_tone}, \textbf{parasympathetic\_tone}
    \item Realistic response to exercise, stress, sleep
\end{itemize}

\subsubsection{Pharmacological Model}
\begin{itemize}
    \item Beta-blockers: Reduce HR, prolong PR
    \item Calcium channel blockers: AV nodal delay
    \item Antiarrhythmics: Class I (Na channel), Class II (beta), Class III (K channel)
    \item Simulate drug effects on conduction and refractoriness
\end{itemize}

\section{Extending the Code}

\subsection{Adding a Custom Pathology Plugin}

\subsubsection{Step 1: Create Plugin Class}
\begin{lstlisting}
from plugins.base import PluginBase

class MyCustomArrhythmia(PluginBase):
    name = "MyCustomArrhythmia"
    description = "My custom ECG pattern"
    
    def __init__(self, params):
        self.params = params
        self.enabled = False
    
    def apply(self, activation_times, conductances):
        if not self.enabled:
            return activation_times, conductances
        
        # Modify activation times or conductances here
        # Example: block AV node
        conductances['av'] *= 0.5
        
        return activation_times, conductances
\end{lstlisting}

\subsubsection{Step 2: Register Plugin}
\begin{lstlisting}
# In main.py or plugin manager
from plugins.my_custom import MyCustomArrhythmia

engine.register_plugin(MyCustomArrhythmia(params))
\end{lstlisting}

\subsubsection{Step 3: Add GUI Controls}
\begin{enumerate}
    \item Add group box to parameter panel
    \item Add enable/disable checkbox
    \item Add parameter sliders (e.g., severity, timing)
    \item Connect to plugin.apply() method
\end{enumerate}

\subsection{Modifying Conduction Geometry}

To change the anatomical connectivity (DAG in Phase 1/2):

\begin{enumerate}
    \item Edit cardiac\_sim/core/tissue/graph.py
    \item Modify adjacency matrix or edge list
    \item Update conduction delays (path\_delay dict)
    \item Test with BFS algorithm
\end{enumerate}

\subsection{Adding New ECG Leads}

\begin{enumerate}
    \item Define new unit vector in cardiac coordinates
    \item Add to LEAD\_MATRIX in lead\_field.py
    \item Update lead names list
    \item GUI automatically displays new lead
\end{enumerate}

\section{Known Limitations}

\begin{enumerate}
    \item \textbf{Single-dipole ECG model}: Limited accuracy for far-field effects
    \item \textbf{No T-wave alternans}: Current model lacks beat-to-beat APD variation
    \item \textbf{No ischemia simulation}: No regional metabolic effects on conduction
    \item \textbf{Simplified atrial geometry}: 8-cell cable vs realistic 3-D atria
    \item \textbf{No autonomic integration}: No HR variability from breathing, stress
    \item \textbf{FHN cell model}: APD ~540 ms (longer than clinical 200--300 ms for ventricle)
\end{enumerate}

\section{Quick Reference: Parameter Summary}

\begin{tabularx}{\textwidth}{|l|c|c|}
\hline
\textbf{Parameter} & \textbf{Range} & \textbf{Default} \\
\hline
HR (bpm) & 30--220 & 70 \\
cycle\_length (s) & 0.27--2.0 & 0.857 \\
av\_delay (s) & 0.05--0.30 & 0.15 \\
av\_conductance & 0.0--1.0 & 1.0 \\
vent\_abs\_refr (s) & 0.25--0.45 & 0.35 \\
amplitude\_scale & 0.1--5.0 & 1.0 \\
\hline
\end{tabularx}

\section{Contact \& Support}

For bug reports, feature requests, or documentation clarifications:

\begin{itemize}
    \item \textbf{Repository}: [GitHub URL to ECG simulator]
    \item \textbf{Documentation}: See Part 1--7 for comprehensive reference
    \item \textbf{Issues}: File issue on GitHub with reproducible steps
    \item \textbf{Discussion}: Review Phase 4+ planning in this Part 8
\end{itemize}

\end{document}
